% !TEX root = main.tex
\section{Homomorphic Densities on Source Value Set}~\label{sec:homomorphic-densities-on-source-value-set}

Since we have a formula of $\mathbb{F}_{n}(b)$ and know the primitive source's distribution, we can build a kind of "measure" to discribe the density of it in $\mathbb{N}$.

In this section, we will always suppose arguments follows \cref{prop:source-set-dijoint-primitive-source}.
\begin{theorem}~\label{thm:infinite-layers}
	For $b\not\in\mathbb{S}$, we have $\forall N > 0$, $\mathbb{F}(b) - \bigcup_{n=0}^{N}\mathbb{F}_{n}(b) \ne \varnothing$.
\end{theorem}
\begin{proof}
	Divide into two situations:\newline
	If $b \not\rightsquigarrow b$. Since $b\not\in\mathbb{S}$, from \cref{prop:source-set-dijoint-primitive-source}, we have $\mathbb{F}_{N}(b) \not\subset \mathbb{S}$. That means $\mathbb{F}_{N + 1}(b) \ne \varnothing$.
	From \cref{prop:non-cycle-layer-independent}, $\forall N \ge 0$ there is
	\[\mathbb{F}(b) - \bigcup_{n=0}^{N}\mathbb{F}_{n}(b) \supset \mathbb{F}_{N+1}(b) - \bigcup_{n=0}^{N}\mathbb{F}_{n}(b) = \mathbb{F}_{N+1}(b) \ne \varnothing.\]
	If $b \rightsquigarrow b$. Suppose it's a $t$-period cycle, then $b \in\mathbb{F}_{t}(b)$. From \cref{thm:layers-of-reduce-space-and-source-value-set}, we can just pick $a = \beta^{\Phi_1}\cdot b + \frac{\gamma}{\alpha}\cdot(\beta^{\Phi_1} - 1)\in\mathbb{F}_{t}(b)$, we have $a \rightsquigarrow b \not\rightsquigarrow a$. From previous discussion, $\mathbb{F}(a)$ can not be covered within finity layers, thus the $\mathbb{F}(b) \supset \mathbb{F}(a)$ also can't be covered within finity layers.
\end{proof}

\cref{thm:infinite-layers} demonstrates that we can never resolve the Collatz conjecture through any method of case analysis, even if we partition the problem into various infinite sets. It's because different layers are in defferent patterns. Though we can tell any given layer's pattern, that's not enough at all. Hence, we shall instead attempt to identify certain homomorphisms in the hope of capturing specific measure-theoretic properties that, we trust, may ultimately aid us in unraveling this conjecture.

From Tao's result\cite{tao2022orbitscollatzmapattain}, when $(\alpha, \beta, \gamma) = (3, 2, 1)$, we have the logarithm density of $\mathbb{F}(1)$ is $\frac{1}{2}$. Suppose there's another number $b$ that $b \not\rightsquigarrow 1$, from \cref{prop:source-set-disjoint} we have $\mathbb{F}(b)\cap\mathbb{F}(1) = \varnothing$, that means the logarithm density of $\mathbb{F}(b)$ will be $0$. Since $\mathbb{F}(b)$ can't be covered by finity layers. Thus, if there a homomorphism that keeps the logarithm measure, or keeps the logarithm measure by layer count, then we can prove that $b$ will not exists. That means no one will not fall into cycle can escape from $1$. However, although Tao's results are classic and very close to the original conjecture, we cannot use them directly; instead, we must prove it really keeps the homomorphism at first.

Summary, if we can find a such homomorphic density and prove the density of $\mathbb{F}(1)$ is almost $1/2$, we then prove \cref{conj:source-value-of-1-cover-all-odd-numbers}, which derives to the classic Collatz conjecture. 
At the same time, we have to say, under the general condition of \cref{prop:source-set-dijoint-primitive-source} instead of $(3,2,1)$, there may not be a conclusion like the classic conjecture, but we can still classify their source value sets.

Now let's introduce a kind of density:
% \begin{definition}~\label{def:arithmetic-density}
% 	Given $N\in\mathbb{N}$ large enough. Define $\mu: \mathcal{P}(\mathbb{N}_{\ge 1}) \to [0, 1]$ as
% 	\[ \mu_{N} A = \frac{\card\{A \le N\}}{N} \]
% 	called $N$-finite arithmetic density. And without causing ambiguity, we also abbreviate it as $\mu A$.
% \end{definition}
% \begin{proposition}
% 	 $\mu\mathbb{N} = 1$, $\mu\varnothing = 0$.
% \end{proposition}
% \begin{proposition}
% 	If $A \cap B = \varnothing$, then $\mu(A\cup B) = \mu A + \mu B$.
% \end{proposition}
% \begin{proposition}
% 	$\mu(\mathbb{N} - A) = 1 - \mu{A}$.
% \end{proposition}
% \begin{proposition}
% 	Suppose $A$ is an infinite set, given $a \in \mathbb{N}$, $\mu(aA) = \frac{1}{a}\mu{A}$.
% \end{proposition}
% \begin{proposition}
% 	$\mu(\beta\mathbb{N}) = \frac{1}{\beta}$, thus $\mu\widehat{\mathbb{H}}_{{\beta}} = \frac{\beta - 1}{\beta}$.
% \end{proposition}
% \begin{proposition}
% 	Given $b \in \mathbb{N}$ that $b\ll N$, we have
% 	\[ \mu(A + b) = \mu{A} + \mathcal{O}(b/N) \approx \mu{A}. \]
% \end{proposition}
\begin{definition}
 	Given $A, B \subset \mathbb{N}$ are infinite set, define
	\[ \mu'(A, B) = \lim_{N\to \infty}\frac{\card\{A \le N\}}{\card\{B \le N\}} \]
	called the relative density between $A$ and $B$.
\end{definition}
From this point onward, unless otherwise specified, we shall always assume sets are infinite.
\begin{theorem}
	Given $a, b\in\widehat{\mathbb{N}}_{\beta} - \mathbb{S}$, we have $\mu'(\mathbb{F}_{1}(a), \mathbb{F}_{1}(b)) = 1$.
\end{theorem}
\begin{proof}
	\begin{align*}
		&\mu'(\mathbb{F}_{1}(a), \mathbb{F}_{1}(b)) \\
		=&\lim_{N\to\infty}\frac{\card\{\mathbb{F}_{1}(a) \le N\}}{\card\{\mathbb{F}_{1}(b) \le N\}} \\
		=&\lim_{N\to\infty}\frac{\card\{k\in\mathbb{N} \mid \frac{1}{\alpha}\beta^{k\Phi_1}\ker a - \frac{\gamma}{\alpha} \le N\}}{\card\{k\in\mathbb{N} \mid \frac{1}{\alpha}\beta^{k\Phi_1}\ker b - \frac{\gamma}{\alpha} \le N\}} \\
		=&\lim_{N\to\infty}\frac{\card\{k\in\mathbb{N} \mid k \le \frac{1}{\Phi_1} \log_{\beta}(\alpha N + \gamma) - \frac{1}{\Phi_1}\log_{\beta}\ker a\}}{\card\{k\in\mathbb{N} \mid k \le \frac{1}{\Phi_1} \log_{\beta}(\alpha N + \gamma) - \frac{1}{\Phi_1}\log_{\beta}\ker b\}} \\
		=&\lim_{N\to\infty}\frac{\lfloor\frac{1}{\Phi_1} \log_{\beta}(\alpha N + \gamma) - \frac{1}{\Phi_1}\log_{\beta}\ker a\rfloor}{\lfloor \frac{1}{\Phi_1} \log_{\beta}(\alpha N + \gamma) - \frac{1}{\Phi_1}\log_{\beta}\ker b\rfloor} \\
		=&\lim_{N\to\infty}\frac{\frac{1}{\Phi_1} \log_{\beta}(\alpha N + \gamma)}{\frac{1}{\Phi_1} \log_{\beta}(\alpha N + \gamma)} \\
		=&1. 
	\end{align*}
\end{proof}
% \begin{remark}
% 	We have to note that $\mu A$ is always a shorthand for $\mu(A \cap \mathbb{H}_{\ge 1})$. Rule 5 can be derived from superfinite additivity using Loeb's construction, but we list it here directly as part of the definition without proof. We can derive $\mu \mathbb{N} = 0$, but this does not violate the fact that the density of $\mathbb{N}$ is $1$ in standard analysis, because $\mathbb{N}$ in standard analysis corresponds to $\nonstd{\mathbb{N}}$.
% \end{remark}
% \begin{example}
% 	Define $\omega(a) = \frac{1}{a\cdot\ln H}$. Given $A\subset \mathbb{H}_{\ge 1}$, we define
% 	\[ \delta A = \sum_{a\in A}\omega(a) = \frac{1}{\ln H}\sum_{a\in A}^{a\in \mathbb{H}_{\ge 1}}\frac{1}{a}. \]
% 	We can tell $\delta$ is the hyperfinite logarithm density. And we have
% 	\begin{itemize}
% 		\item[1.] $\delta\mathbb{H}_{\ge 1} = \frac{\ln H + \gamma_{\mathrm{e}}}{\ln H}$, here $\gamma_{\mathrm{e}}$ is the Euler constant. We have $\delta\varnothing = 0$ and $\forall a\in \mathbb{H}_{\ge 1}$, $\delta\{a\} = \st\omega(a) = 0$.
% 		\item[2.] If $A, B \in \mathcal{L}(\mathbb{H}_{\ge 1})$ and $A \cap B = \varnothing$, then
% 		\[ \delta(A\cup B) = \frac{1}{\ln H}\sum_{a\in A\cup B}^{a\in \mathbb{H}_{\ge 1}}\frac{1}{a} = \frac{1}{\ln H}\sum_{a\in A}^{a\in \mathbb{H}_{\ge 1}}\frac{1}{a} + \frac{1}{\ln H}\sum_{a\in B}^{a\in \mathbb{H}_{\ge 1}}\frac{1}{a} = \delta A + \delta B. \]
% 		\item[3.] Suppose $A \in \mathcal{L}(\mathbb{H}_{\ge 1})$, given $a \in \mathbb{N}$, we have
% 		\begin{align*}
% 			\delta(aA \cap \mathbb{H}_{\ge 1}) &= \frac{1}{\ln H}\sum_{s\in A}^{a\cdot s\in \mathbb{H}_{\ge 1}}\frac{1}{a\cdot s} \\
% 			&= \frac{1}{a} \frac{1}{\ln H}\sum_{s\in A}^{s\in \mathbb{H}_{\ge 1}/a}\frac{1}{s} \\
% 			&= \frac{1}{a}\left(\delta{A} - \frac{1}{\ln H}\sum_{H/a < s \le H}^{a\in A}\frac{1}{s}\right).
% 		\end{align*}
% 		Notice $0 \le \frac{1}{\ln H}\sum_{H/a < s \le H}^{a\in A}\frac{1}{s}\le \frac{\ln H - \ln (H/a)}{\ln H} = \frac{\ln a}{\ln H} \approx 0$, thus we have $\delta(aA \cap \mathbb{H}_{\ge 1}) = \frac{1}{a}\delta{A}$.

% 		\item[4.] Suppose $A \in \mathcal{L}(\mathbb{H}_{\ge 1})$, given $b \in \mathbb{N}$, simillarly, we have
% 		\begin{align*}
% 			&\left\vert\delta((A+b) \cap \mathbb{H}_{\ge 1}) - \delta(A\cap \mathbb{H}_{\ge 1})\right\vert \\
% 			=&\frac{1}{\ln H}\left\vert\sum_{a\in A}^{a\le H}\frac{1}{a} - \sum_{a\in A}^{a+b \le H}\frac{1}{a + b}\right\vert \\
% 			=&\frac{1}{\ln H}\left\vert\sum_{a\in A}^{a\le H}\left(\frac{1}{a}-\frac{1}{a + b}\right) + \sum_{a\in A}^{H-b < a \le H}\frac{1}{a + b}\right\vert \\
% 			=&\frac{1}{\ln H}\left\vert\sum_{a\in A}^{a\le H}\frac{b}{a(a + b)} + \sum_{a\in A}^{H-b < a \le H}\frac{1}{a + b}\right\vert \\
% 			\le&\frac{1}{\ln H}\left\vert\frac{b\pi^{2}}{6} + \ln (H+b) - \ln(H)\right\vert \\
% 			=&0.
% 		\end{align*}
% 		Thus $\delta((A+b) \cap \mathbb{H}_{\ge 1}) = \delta(A\cap \mathbb{H}_{\ge 1})$.
% 		\item[5.] Since $A_n \in \mathcal{L}(\mathbb{H}_{\ge 1})$, then $A = \bigcup_{n=1}^{\infty}A_n \in \mathcal{L}(\mathbb{H}_{\ge 1})$, thus $\delta A$ exists and $\delta A = \sum_{n=0}^{\infty}\delta A_n$.
% 	\end{itemize}
% 	Sum up we proved $\delta$ is an non-standard arithmetic density.
% \end{example}
\begin{theorem}~\label{thm:proportional-of-density-of-source-value-set}
	$\forall a, b \in \widehat{\mathbb{N}}_\beta - \mathbb{S}$ that $a\not\rightsquigarrow b$ and $b\not\rightsquigarrow a$, there is
	\[\mu'(\mathbb{F}_{n}(b), \mathbb{F}_{n}(a)) = \frac{B(a;n)}{B(b;n)}. \]
	% Here $B(t;n) = \prod_{k=0}^{n-1}\sum_{i=1}^{\alpha-1}\beta^{v(i;t_{-k})}$ only depends on the distribution of source value of $t$ modulo $\alpha$.
\end{theorem}
\begin{lemma}
	$\card\{\mathbb{F}_{n}(a) \le N\} = \card\{\bm{v} \in \mathbb{V}_{n}(a) \mid Z_{n}(N;\bm{v}) \ge 1 \}$. Actually, we are more conserned about the count of $\simp\bm{v}$.
\end{lemma}
\begin{lemma}~\label{lem:finite-layers-le-M}
	Suppose $b\not\rightsquigarrow b$.
	$\forall M > 0$, $\exists H > 0$ depends on $M$ such that $\forall n \ge H$, there is $\card\{ \mathbb{F}_{n}(b)\le M \} = 0$.
\end{lemma}
\begin{proof}
	Since $b\not\rightsquigarrow b$, thus $\forall i, j \ge 0$, $\mathbb{F}_{i}(b)\cap \mathbb{F}_{j}(b) = \varnothing$.
	We can construct a sequence of set that $A_{n} = \{x\in\mathbb{N} \mid x \le M\} - \bigcup_{k=0}^{n-1}\mathbb{F}_{k}(b)$. We have $A_{n+1} \subset A_{n}$. Thus $\{A_n\}_{n=1}^{\infty}$ has a minimal element. Suppose it's $A_m$. Then let $ H = \min\{k\in\mathbb{N}\mid A_{k} = A_{m}\}$, $\forall n\ge H$, we have $\card\{\mathbb{F}_{n}(b)\le M\} = 0$.

	Mow let's consider the relationship between $H$ and $M$. Since $F_{n}(b)$ is a infinite set, thus given $ M' > M$, if $\forall n > H(M)$, there is $\{M < F_{n}(b) \le M' \} = \varnothing$ that $\card\{\mathbb{F}_{n}(b)\le M'\} = \card\{\mathbb{F}_{n}(b)\le M\}$, then $H(M') = H(M)$. If $\exists n\ge H(M)$ such that $\{M < F_{n}(b) \le M' \} \ne \varnothing$, then we have $H(M') > H(M)$. Sum up, we have $H(M') \ge H(M)$, thus $H$ is increasing according to $M$.
\end{proof}
\begin{lemma}~\label{lem:source-value-set-finite-decomposition}
	Suppose $b\not\rightsquigarrow b$. For $H$ in \cref{lem:finite-layers-le-M}, we have
	\[ \mu\mathbb{F}(b) = \mu\bigcup_{n=0}^{H-1}\mathbb{F}_{n}(b) = \sum_{n=0}^{H-1}\mu\mathbb{F}_{n}(b). \]
\end{lemma}
% \begin{lemma}
% 	If $b\not\rightsquigarrow b$, then $\forall n\in\mathbb{N}$, we have
% 	\[ \frac{\mu\mathbb{F}_{n+1}(b)}{\mu\ker\mathbb{F}_{n}(b)} = \alpha\frac{(\beta^{H(n)} - 1)\beta^{\Phi_1}}{\beta^{H(n)}(\beta^{\Phi_1} - 1)} = \frac{\alpha\beta^{\Phi_1}}{\beta^{\Phi_1} - 1}\left(1-\frac{1}{\beta^{H(n)\Phi_1}}\right). \]
% 	Here $H(n)$ keeps all $\mathbb{F}_{n}(b)$ less than given limit.
% \end{lemma}
% \begin{proof}
% 	Since $b\not\rightsquigarrow b$, we have $\forall n \in \mathbb{N}$, $\mathbb{F}_{n+1}(b)\cap\ker\mathbb{F}_{n}(b) = \varnothing$.
% 	Then from \cref{cor:source-value-set-equivlants-its-pre-ker-set}, we have $\alpha\mathbb{F}_{n+1}(b) + \gamma = \mathbb{B}\otimes \ker\mathbb{F}_{n}(b)$. Thus we have
% 	\begin{align*}
% 		\mu(\alpha\mathbb{F}_{n+1}(b) + \gamma) &= \mu\mathbb{B}\otimes \ker\mathbb{F}_{n}(b) \\
% 		&= \frac{1}{N}\sum_{x\in\mathbb{B}}^{x\le N} \card\left\{\ker\mathbb{F}_{n}(b) \le \frac{N}{x} \right\} \\
% 		&= \frac{1}{N}\sum_{k=0}^{\frac{1}{\Phi_1}\log_{\beta}N} \card\left\{\ker\mathbb{F}_{n}(b) \le \frac{N}{\beta^{k\Phi_1}} \right\} \\
% 		&= \frac{1}{N}\sum_{k=0}^{\frac{1}{\Phi_1}\log_{\beta}N}\sum_{i=1}^{\alpha-1} \card\left\{\widetilde{\mathbb{F}_{n}(b)}_i \le \frac{N}{\beta^{k\Phi_1}}\big/\beta^{v(i)} \right\} \\
% 		&= \frac{1}{N}\sum_{k=0}^{\frac{1}{\Phi_1}\log_{\beta}N}\sum_{i=1}^{\alpha-1} \beta^{-v(i)}\card\left\{\widetilde{\mathbb{F}_{n}(b)}_i \le \frac{N}{\beta^{k\Phi_1}} \right\} \\
% 		&= \frac{1}{N}\sum_{k=0}^{\frac{1}{\Phi_1}\log_{\beta}N}\sum_{i=1}^{\alpha-1} \beta^{-v(i)}\card\left\{\widetilde{\mathbb{F}_{n}(b)}_i \le \frac{N}{\beta^{k\Phi_1}} \right\}
% 	\end{align*}
% \end{proof}
% \begin{lemma}
% 	If $b\not\rightsquigarrow b$, $\forall n\ge 1$, we have
% 	\[ k(n) = \frac{\mu\mathbb{F}_{n+1}(b)}{\mu\mathbb{F}_{n}(b)} = \eta \cdot \left(1-\frac{1}{\beta^{H(n)\Phi_1}}\right), \]
% 	here $\eta = \frac{\alpha\beta^{\alpha\Phi_1}}{\beta^{\alpha\Phi_1} - 1}\sum_{i=1}^{\alpha-1}\beta^{-v(i)-i\Phi_1}$ is a constant.
% \end{lemma}
% \begin{proof}[Proof of \cref{thm:proportional-of-density-of-source-value-set}]
% 	\[ \frac{\mu\mathbb{F}(b)}{\mu \mathbb{B}} = \frac{\alpha}{\ker b}\cdot\frac{\mu\mathbb{F}(b)}{\mu\mathbb{F}_{1}(b)} = \frac{\alpha}{\ker b}\cdot\sum_{n=0}^{\infty}\prod_{m=0}^{n}k(m). \]
% 	Denote $K = \sum_{n=0}^{\infty}\prod_{m=0}^{n}k(m)$ is a constant. Given $a, b\in\widehat{\mathbb{N}}_{\beta} - \mathbb{S}$ that $a\not\rightsquigarrow a$ and $b\not\rightsquigarrow b$, we can derive that $\frac{\mu\mathbb{F}(a)}{\mu\mathbb{F}(b)} = \frac{\ker b}{\ker a}$, which exactly is
% 	\[ \ker a \cdot\mu\mathbb{F}(a) = \ker b \cdot\mu\mathbb{F}(b). \]
% \end{proof}

% Now let's consider $K\mu\mathbb{B}$. It radically determines the density of $\mathbb{F}(b)$ and thus the distribution of this Collatz system. $\mu \mathbb{B}$ is a infinitismal, and if $k(m) \ge 1$, then $K$ goes to infinity, thus $K\mu \mathbb{B}$ is indeterminate form, we may guass for certain arguments $(\alpha, \beta, \gamma)$, it is finite and non-zero.
\begin{corollary}
	$\forall a, b \in \widehat{\mathbb{N}}_\beta - \mathbb{S}$, we have $\mu'(\mathbb{F}(b), \mathbb{F}(a)) > 0$.
\end{corollary}
\begin{corollary}
	If $\exists a \in \widehat{\mathbb{N}}_\beta - \mathbb{S}$ such that $\lim_{N\to\infty}\mu\mathbb{F}(a) \ne 0$ , then
	$\forall b\in\widehat{\mathbb{N}}_\beta - \mathbb{S}$ we have $\lim_{N\to\infty}\mu\mathbb{F}(b) \ne 0$.
\end{corollary}
We have to recall that, all conclusions of $\mathbb{F}(b)$ in this section are base on the condition of \cref{prop:source-set-dijoint-primitive-source}.
From \cref{thm:count-function-from-dirichlet-series} we have
\begin{corollary}
    For large enough $c > 0$, we have
    \[ \widetilde{\pi}(\alpha\mathbb{F}_{n+1}(b) + \gamma;x) = \frac{1}{2\pi\rmi}\int_{c-\rmi\infty}^{c+\rmi\infty}\frac{D(\ker\mathbb{F}_{n}(b);s)}{1- \beta^{s\Phi_1}}\cdot\frac{x^{s}}{s}\rmd{s}. \]
\end{corollary}
\begin{proof}
	From \cref{cor:source-value-set-equivlants-its-pre-ker-set}, we have 
	\begin{align*}
		D(\alpha\mathbb{F}_{n+1}(b) + \gamma; s) &= D(\mathbb{B}\otimes \ker\mathbb{F}_{n}(b); s) \\
		&= D(\mathbb{B};s)\cdot D(\ker\mathbb{F}_{n}(b);s) \\
		&= \sum_{k=0}^{\infty}\beta^{ks\Phi}\cdot D(\ker\mathbb{F}_{n}(b);s) \\
		&= \frac{D(\ker\mathbb{F}_{n}(b);s)}{1 - \beta^{s\Phi}}.
	\end{align*}
	Thus from \cref{thm:count-function-from-dirichlet-series}, we have
	\begin{align*}
		\widetilde{\pi}(\alpha\mathbb{F}_{n+1}(b) + \gamma;x) &= \frac{1}{2\pi\rmi}\int_{c-\rmi\infty}^{c+\rmi\infty}D(\alpha\mathbb{F}_{n+1}(b) + \gamma; s)\cdot\frac{x^{s}}{s}\rmd{s} \\
		&= \frac{1}{2\pi\rmi}\int_{c-\rmi\infty}^{c+\rmi\infty}\frac{D(\ker\mathbb{F}_{n}(b);s)}{1- \beta^{s\Phi_1}}\cdot\frac{x^{s}}{s}\rmd{s}.
	\end{align*}
\end{proof}

